\documentclass[10pt,a4paper]{amsart}
\usepackage[margin=19mm]{geometry}
\usepackage{amsmath,amssymb,mathtools}
\usepackage[scr=rsfs,cal=euler]{mathalfa}
\usepackage{xfrac}
\usepackage[unicode,hidelinks,bookmarks=false]{hyperref}
\newcommand{\ie}{i.e.\ }
\newcommand{\cf}{cf.\ }
\newcommand{\resp}{respectively\ }
\newcommand{\Subsection}{\subsection}
\providecommand{\nobreakdash}{\nobreak\hbox{-}}
\setlength{\emergencystretch}{2em}
\multlinegap0pt
\usepackage{tikz}
\newtheorem{theorem}{Theorem}
\title{Gaussian Behavior and Geometric Gaps in Decompositions from Recurrences with Zero Coefficients}
\author{Sajad Salami}
\date{Source: arXiv:2604.15447v1 (CC BY 4.0)}
\begin{document}
\maketitle

\noindent\emph{Source and licence.} Gaussian Behavior and Geometric Gaps in Decompositions from Recurrences with Zero Coefficients. Version arXiv:2604.15447v1 (CC BY 4.0), distributed by arXiv under the Creative Commons Attribution 4.0 International licence, http://creativecommons.org/licenses/by/4.0/, as recorded in the arXiv licence metadata for this exact version and shown on the abstract page https://arxiv.org/abs/2604.15447v1. Full licence notice: \texttt{licenses/arxiv-2604.15447-CC-BY-4.0.txt}.\par\smallskip

\noindent\textit{Editorial scope.} Counted unit: the article's theorem conj:num\_decomps (source0.tex lines 596--599). The article's complete original proof of it is delivered with it (source0.tex lines 600--625, one \texttt{proof} environment of 26 lines). No mathematics is written, repaired or re-derived: the statement and the proof are the article's own lines, in the article's own notation, and no retained line is substituted or re-flowed.  No further source-local context is needed: every cross-reference of every retained line resolves inside the two delivered spans.  Nothing was omitted from any retained span, so the package carries no omission record.  No retained line contains a citation, so the wrapper carries no bibliography and no bibliography entry.  No retained line contains a figure environment, an image-inclusion command or drawing code, so no image is delivered and the package declares no asset dependency.  Editorial formatting only, in addition: an AMS \texttt{amsart} preamble that restates the article's own declarations and macros used by the retained lines, namely the environments \texttt{theorem} on separate counters, together with the standard packages those lines need.  The retained environments print in this document as Theorem 1 in order of appearance, whereas the article prints the same items with the numbering of its own full document.  The complete native source of the article as distributed in the arXiv e-print for this exact version is bundled unchanged as \texttt{sources/198606/source0.tex}.
	\begin{theorem}
		\label{conj:num_decomps}
		Given a positive integer $N$, let $d(N)$ be the number of distinct legal decompositions of $N$ in terms of the Lagonacci sequence $\{Z_n\}_{n\geq 0}$. The total number of legal binary strings of length $N$, representing the sum of decompositions for all integers up to $Z_{N}$ and denoted by $D(Z_{N})=\sum_{k=0}^{Z_{N}-1}d(k)$, grows asymptotically as $D(Z_{N})\sim C\cdot 2^{N}$ for some constant $C$. Consequently, the number of representations grows exponentially faster than the number of integers, with the average number of decompositions as  given below, $$ \bar{d}_N = \frac{D(Z_N)}{Z_N} \sim (1.509755)^N.$$
	\end{theorem}

	\begin{proof}
		We use the transfer matrix method. A legal Lagonacci decomposition is a sum $\sum \epsilon_i Z_i$ where $\epsilon_i \in \{0,1\}$. The rule defining a legal decomposition is that no sum can be simplified using the recurrence relation. For the Lagonacci sequence, $Z_{n+1} = Z_{n-1} + Z_{n-2}$, this means we cannot have a string of coefficients such that $\epsilon_{n-1}=1$ and $\epsilon_{n-2}=1$ while $\epsilon_{n+1}=0$ (and intermediate coefficients are zero), as $Z_{n-1} + Z_{n-2}$ would be replaced by $Z_{n+1}$. This leads to a forbidden substring of coefficients $(\dots, \epsilon_{n-2}, \epsilon_{n-1}, \epsilon_n, \epsilon_{n+1}, \dots) = (\dots, 1, 1, 0, 0, \dots)$, where indices increase to the right.
		
		To count the number of binary strings of length $N$ that avoid the substring `1100`, we define states based on the last three coefficients seen, i.e., $(\epsilon_{i-2}, \epsilon_{i-1}, \epsilon_i)$. The next coefficient, $\epsilon_{i+1}$, determines the transition to a new state $(\epsilon_{i-1}, \epsilon_i, \epsilon_{i+1})$. The transition from state $(1,1,0)$ to $(1,0,0)$ is forbidden, as this would complete the `1100` pattern. All other transitions are allowed. This leads to an $8 \times 8$ transfer matrix $T$, where rows index the starting state and columns index the ending state.
		\[ T = \begin{pmatrix}
			% 000 001 010 011 100 101 110 111 (To)
			1 & 1 & 0 & 0 & 0 & 0 & 0 & 0 \\ % From 000
			0 & 0 & 1 & 1 & 0 & 0 & 0 & 0 \\ % From 001
			0 & 0 & 0 & 0 & 1 & 1 & 0 & 0 \\ % From 010
			0 & 0 & 0 & 0 & 0 & 0 & 1 & 1 \\ % From 011
			1 & 1 & 0 & 0 & 0 & 0 & 0 & 0 \\ % From 100
			0 & 0 & 1 & 1 & 0 & 0 & 0 & 0 \\ % From 101
			0 & 0 & 0 & 0 & 0 & 1 & 0 & 1 \\ % From 110 (100 is forbidden)
			0 & 0 & 0 & 0 & 0 & 0 & 1 & 1   % From 111
		\end{pmatrix}. \]
		Applying the characteristic equation to the matrix $T$ yields:
		\[ Q(x) = \det(xI - T) = x^4(x - 2)(x^3 - x - 1) \]
		%	
		The roots of this polynomial provide the growth rates for the system. The factor $P(x)=x^3 - x - 1$ corresponds to the Lagonacci sequence's own growth, while the factor $x - 2$ identifies the Perron-Frobenius eigenvalue
		\( \alpha = 2 \),  indicating that the number of legal binary strings of length $N$ grows as $D(Z_N) \sim C \cdot 2^N$,
		for some constant C.
		We now compare this to the growth of the underlying sequence $Z_N$, which is governed by the plastic ratio $\lambda_1 \sim 1.324718$. 
		Since $\alpha = 2 > \lambda_1$, the number of representations grows exponentially faster than the number of integers. Specifically, the average number of decompositions for a typical integer of size $\sim Z_N$ grows at the rate:
		\[ \bar{d}_N = \frac{D(Z_N)}{Z_N} \sim \left( \frac{2}{\lambda_1} \right)^N \sim (1.509755)^N. \]
		
	\end{proof}

\end{document}
